\section{Introduction}\label{sec:intro}

A graph is a finite set of points, called vertices, together with lines, called edges, each joining two vertices. The adjacency matrix of a graph is a table of zeros and ones that records which pairs of vertices are joined. This note takes a graph with $N$ vertices and adjacency matrix $A$ (defined in \cref{sec:defs}) and answers three questions, using only matrix multiplication and induction.
\begin{enumerate}
\item[(Q1)] What does the entry $(A^k)_{ij}$ of the $k$-th power of $A$ count? \emph{Answer:} the number of walks of length $k$ from vertex $v_i$ to vertex $v_j$; walks may repeat vertices, so this is not a count of paths (\cref{thm:counting}).
\item[(Q2)] What is a path between two vertices? \emph{Answer:} a walk between them in which no vertex occurs twice (\cref{def:path}).
\item[(Q3)] How can one decide from $A$ alone whether the graph is connected? \emph{Answer:} for $N\ge 2$, the graph is connected if and only if every entry off the main diagonal of $A+A^2+\dots+A^{N-1}$ is positive (\cref{thm:sum}).
\end{enumerate}
The questions matter because the answer to Q3 turns a question about a picture into a computation on a table of numbers. The answers are known; the walk-counting statement appears in several sets of notes \cite{galwaynotes,gotti18211,bai2025adjacency}, the lemma that a walk contains a path in a graph theory text \cite{grinberg2025graph}, and the sum test in a journal article \cite{paramadevan2021adjacency}. Their proofs are written out in full here, because the sources sketch or omit some steps. The support for the answers is these proofs and the examples in \cref{sec:example}, which can be checked by hand; the limits of the test are stated at the end of that section. \Cref{sec:defs} gives the definitions, including the answer to Q2; \cref{sec:counting} answers Q1; \cref{sec:test} answers Q3, which needs the other two.
